\small
\noindent\textbf{Resumen}
%Abstract below

La enfermedad renal crónica (ERC) es una pérdida progresiva de la función renal que afectó a más de 800 millones de personas en 2017, 
con frecuencia subdiagnosticada por presentar síntomas en etapas avanzadas. Su detección temprana podría reducir significativamente su 
carga clínica y económica. El siguiente estudio empleó un enfoque cuantitativo no experimental basado en redes bayesianas gaussianas (GBN) 
para modelar las relaciones de dependencia probabilística entre biomarcadores clínicos asociados a la ERC, utilizando la base de datos 
ckd.csv. Se construyeron tres grafos acíclicos dirigidos (DAGs) con distintos subconjuntos de variables, definidos con apoyo de estudiantes 
y profesores de medicina, ajustados con bn.fit() en R mediante el paquete bnlearn y evaluados utilizando los criterios BIC y AIC. La DAG 2 
presentó el mejor desempeño global y fue seleccionada para realizar las consultas probabilísticas propuestas por especialistas. Los 
resultados mostraron probabilidades de 62.77 % de creatinina sérica elevada dado un nivel alto de sodio, 62.84 % al considerar 
simultáneamente sodio y potasio elevados, 43.13 % de hiperpotasemia ante creatinina y urea elevadas, 80.96 % de hemoglobina menor a 12 g/dL 
ante una creatinina de 3.0 mg/dL y un conteo de eritrocitos de 3.5 millones/\(\mu\)L, y 100 % de creatinina elevada ante presión arterial 
superior a 140 dentro del modelo ajustado. Asimismo, el modelo no paramétrico mejoró los valores de BIC y AIC de la DAG 2 respecto al modelo 
lineal, aunque las consultas probabilísticas se realizaron sobre este último. Finalmente, se identificó como limitación la imposibilidad de 
incorporar directamente variables categóricas, como el diagnóstico de ERC, en una red exclusivamente gaussiana, proponiéndose como extensión 
futura el uso de redes bayesianas gaussianas condicionales o redes discretas.
